\section{User Testing Feedback}
\label{sec:user-testing}

Some elements of functional testing were conducted with (6) volunteers, required to test the system's capability of registering and classifying gait-based identities.
All users were required to read a Participant Information Sheet (PIS) and fill out an Informed Consent Form (ICF) before participating in the testing sessions.
Templates of these documents can be seen in Appendix \ref{app:participant-information-sheet} and Appendix \ref{app:informed-consent-form} respectively.

To discern the system's usability and understand the potential for real-world adoption, volunteers were asked to complete a questionnaire following the testing sessions, containing the following questions:

\begin{enumerate}
	\item How easy was it to understand the instructions during the registration/on-boarding process? (Very difficult, Difficult, Neutral, Easy, Very easy)
	\item Did you encounter any technical issues or confusion during onboarding? Please specify if yes. (Yes, No)
	\item How would you rate the speed of the authentication system? (Very slow, Slow, Acceptable, Fast, Very fast)
	\item How convenient was the passive authentication compared to using a physical token, or using facial/fingerprint recognition? (Very inconvenient, Inconvenient, Neutral, Convenient, Very convenient)
	\item Did the authentication process feel seamless and non-intrusive? (Not at all, Slightly, Moderately, Mostly, Completely)
	\item Would you be willing to use this system in your workplace or daily environment? (Definitely not, Probably not, Not sure, Probably, Definitely)
	\item Do you have any privacy, ethical, or security concerns regarding this type of behavioural biometric identification? Please specify if yes. (Yes, No)
\end{enumerate}

The responses from the volunteers are summarised below:

\paragraph{Response 1} Volunteer 1 was a 21 year old male, he claimed to have found the registration process easy, without experiencing any technical issues.
He rated the speed of the authentication system as 'Fast', and found the passive authentication process to be 'Convenient'.
He found the authentication process to be 'Mostly' seamless and non-intrusive, and stated he would 'Probably' be willing to use the system in a workplace, and had no privacy reservations.

\paragraph{Response 2} Volunteer 2 was a 21 year old male, he found the registration process 'Very easy', and did not experience any technical issues.
He rated the speed of the authentication system as 'Acceptable', and found the passive authentication process to have 'Convenient'.
He found the authentication process to be 'Completely' seamless and non-intrusive, stating he would 'Definitely' use the system in a workplace, and had no privacy concerns.

\paragraph{Response 3} Voluteer 3 was a 23 year old female, she found the registration process 'Easy', but stated 'It was unclear how long to wait for the system to respond when waiting for authentication'.
She rated the speed of the authentication system as 'Slow', and found the passive authentication process to be 'Neutral'.
She stated the authentication process was 'Slightly' seamless and non-intrusive, but admitted she would 'Probably not' be willing to use the system in a workplace, and had privacy concerns: Yes --- 'I have concerns about how this type of system could be used for surveillance, the idea of being tracked without knowing makes me uncomfortable'.

\paragraph{Response 4} Volunteer 4 was a 53 year old male, he found the registration process 'Easy', and did not experience any technical issues.
He rated the speed of the authentication system as 'Acceptable', and found the passive authentication process to be 'Neutral'.
He found the authentication process to be 'Moderately' seamless and non-intrusive, but stated he was 'Not sure' if he would be willing to use the system in a workplace, yet had no privacy concerns.

\paragraph{Response 5} Volunteer 5 was a 23 year old female, she found the registration process 'Easy', but pointed out the fact that her silhouette was incomplete due to the low contrast between her trousers and the background.
She rated the speed of the authentication system as 'Acceptable', and rated the passive authentication process as 'Neutral'.
She found the authentication process to be 'Moderately' seamless and non-intrusive, but stated she would 'Probably not' be willing to use the system in a workplace, and had no privacy concerns but mentioned that she: 'Was satisfied with the convenience of using a key card'.

\paragraph{Response 6 (Unseen Test Subject)} Volunteer 6 was a 56 year old female, she found the registration process 'Easy', and did not experience any technical issues.
She rated the speed of the authentication system as 'Acceptable', and rated the passive authentication process as 'Neutral'.
She found the authentication process to be 'Moderately' seamless and non-intrusive, but stated she would 'Probably not' be willing to use the system in a workplace, and had no privacy concerns.

These responses convey a mixed but overall positive outlook on \ac{gr} based authentication, with most volunteers finding the registration process and subsequent authentication process easy, experiencing no technical issues.
However, some did raise concerns about privacy and the potential for surveillance, echoing past sentiments from the survey conducted in Chapter \ref{chap:context}.
Another volunteer simply stated that she was happy with the convenience of \ac{rfid} key cards.
From a usability and performance perspective, most volunteers rated the authentication speed positively and found the passive authentication process to be convenient.

Overall, these responses, when combined with the previous survey results, show that while there is an interest in \ac{gr} based authentication, significant concerns about privacy and data handling remain.
This makes sense, as technology that can authenticate users based on biometrics without their participation, could also be used for more nefarious purposes, such as non-consensual surveillance.

